\lysh{15689}{4}
\begin{alist}
\item Η $f$ είναι άρτια, οπότε η γραφική της παράσταση είναι συμμετρική ως προς τον άξονα $y'y$. Θεωρώντας το συμμετρικό του δοσμένου σχήματος ως προς τον $y'y$ προκύπτει η γραφική παράσταση του διπλανού σχήματος στο οποίο φαίνεται η γραφική παράσταση της $f$ για όλες τις τιμές του $x$.
\begin{center}
\begin{tikzpicture}
\begin{axis}[
    axis lines=middle,
    xmin=-3.2, xmax=3.2,
    ymin=-1.2, ymax=2,
    xtick={-3,-2,-1,1,2,3},
    ytick={-1,1,2},
    grid=both,
    grid style={gray!30},
    thick
]
\addplot[domain=0:2.4, samples=100,maincolor] {2*sqrt(x)-x^2};
\addplot[domain=-2.4:0, samples=100,secondarycolor] {2*sqrt(-x)-(-x)^2};
\addplot[mark=none, draw=black,dashed] coordinates {(1,1) (1,0)};
\addplot[mark=none, draw=black,dashed] coordinates {(-1,1) (-1,0)};
\end{axis}
\end{tikzpicture}
\end{center}
\item Από τη γραφική παράσταση της $f$ προκύπτει ότι η $f$ είναι γνησίως αύξουσα σε καθένα από τα διαστήματα $[-\sqrt{5}, -1]$ και $[0, 1]$ και γνησίως φθίνουσα σε καθένα από τα διαστήματα $[-1, 0]$ και $[1, \sqrt{5}]$. Η ελάχιστη τιμή της $f$ είναι ίση με $0$ και προκύπτει όταν $x = -\sqrt{5}, x = 0$ ή $x = \sqrt{5}$. Η μέγιστη τιμή της είναι ίση με $1$ και προκύπτει όταν $x = -1$ ή $x = 1$.
\item 
\begin{rlist}
\item Είναι γνωστό ότι η συνάρτηση ημίτονο είναι γνησίως αύξουσα και η συνάρτηση συνημίτονο είναι γνησίως φθίνουσα στο διάστημα $\left(0, \dfrac{\pi}{2}\right)$. \\
Επιπλέον $\hm \dfrac{\pi}{4} = \syn \dfrac{\pi}{4} = \dfrac{\sqrt{2}}{2}$, οπότε με $\theta \in \left(\dfrac{\pi}{4}, \dfrac{\pi}{2}\right)$ έχουμε:
\[\theta > \dfrac{\pi}{4} \Rightarrow \hm \theta > \hm \dfrac{\pi}{4} \text{ και } \theta > \dfrac{\pi}{4} \Rightarrow \syn \theta < \syn \dfrac{\pi}{4}, \text{ επομένως}\]
\[\hm \theta > \hm \dfrac{\pi}{4} = \syn \dfrac{\pi}{4} > \syn \theta, \text{ άρα } \hm \theta > \syn \theta\]
\item Με τη βοήθεια του ερωτήματος γ)i. και δεδομένου ότι οι αριθμοί $\hm \theta$ και $\syn \theta$ περιέχονται στο διάστημα $[0, 1]$, όπου η $f$ είναι γνησίως αύξουσα, έχουμε:
\[\hm \theta > \syn \theta \Rightarrow f(\hm \theta) > f(\syn \theta)\]
\end{rlist}
\textbf{Εναλλακτική αντιμετώπιση του γ) i.} \\
Η ανισοτική σχέση ανάμεσα στα $\hm \theta$ και $\syn \theta$ στο διάστημα $\left(\dfrac{\pi}{4}, \dfrac{\pi}{2}\right)$ μπορούσε να προκύψει επίσης και από τον τριγωνομετρικό κύκλο ή τις γραφικές παραστάσεις των αντίστοιχων συναρτήσεων.
\end{alist}

